\documentclass[conference]{IEEEtran}
\IEEEoverridecommandlockouts
\usepackage{amsmath,amssymb,amsfonts}
\usepackage{textcomp}
\usepackage[T1]{fontenc}
\usepackage[utf8]{inputenc}
\usepackage{url}
\usepackage{xspace}

% --- spacing of section titles kept default to preserve IEEEtran layout ---

\newcommand{\beam}{BEAM\xspace}
\newcommand{\otp}{OTP\xspace}
\newcommand{\dt}{DT\xspace}

\begin{document}

\title{Twins as Processes: A Pure Erlang/OTP\\ Framework for Digital Twins Across the\\ Edge--Cloud Continuum}

\author{\IEEEauthorblockN{[Author Name]}
\IEEEauthorblockA{\textit{School of Computing Science} \\
\textit{Simon Fraser University}\\
Burnaby, BC, Canada \\
{[email]}}
}

\maketitle

\begin{abstract}
The digital twin has matured from a NASA-era spacecraft analogue into a
foundational construct of Industry~4.0, yet its conceptual and architectural
fragmentation persists: the term spans one-way digital shadows, full
closed-loop cyber-physical systems, and cognitive agents, with no agreement on
the runtime substrate that should host them. Drawing on the design space
codified in the \emph{Handbook of Digital Twins}, this paper argues that the
recurring requirements a twin framework must satisfy---per-asset autonomy,
hierarchical composition from component to system-of-systems, variable
synchronization, fault containment, live model evolution, and deployment that
ranges from an embedded device to a cloud cluster---are precisely the
properties that the actor model provides, and that the Erlang/Open Telecom
Platform (\otp) runtime, the \beam, realizes them more directly than either
vendor twin platforms or actor frameworks ported onto general-purpose runtimes.
We develop the correspondence in detail, mapping each axis of the Handbook
taxonomy onto a native \beam mechanism, and from it derive an actor-oriented
framework in pure Erlang. A twin is a supervised \texttt{gen\_statem} process
that carries an explicit split between observed and modelled state, whose
divergence operationalizes the digital-shadow-to-digital-twin maturity ladder;
asset topology is expressed as a supervision hierarchy of composite twins;
transport is isolated behind adapters; and registry, history, and time are
pluggable behaviours. The same twin module runs as a thin leaf under AtomVM on
a microcontroller and as a composite in a distributed cluster, with live code
replacement updating a running model without halting the asset it mirrors. We
ground the design in a water-distribution case study built on the C-Town
benchmark and a simulator-replay adapter, and discuss the limits of the
approach, notably numerical modelling and ecosystem maturity.
\end{abstract}

\begin{IEEEkeywords}
digital twin, actor model, Erlang, BEAM, OTP, edge computing, cyber-physical
systems, supervision, Industry 4.0, AtomVM
\end{IEEEkeywords}

\section{Introduction}
The digital twin (\dt) has become one of the organizing ideas of digital
transformation, promising a virtual counterpart that is bound to a physical
asset closely enough to monitor, predict, and ultimately steer it throughout
its lifecycle~\cite{handbook,tao2019,fuller2020}. The breadth of that promise
is also the source of a persistent difficulty. As Agrawal and Fischer observe
in their conceptual treatment, practitioners routinely disagree on what a twin
even is: some take it to mean a complex, predictively capable real-time model,
others a lightweight digital representation, and the boundary between a twin and
the adjacent vocabulary of Building Information Modelling, cyber-physical
systems, and Industry~4.0 is blurry in practice~\cite{agrawal2024,fuller2020}.
This ambiguity is not merely terminological. It propagates into architecture,
where it becomes a question of what software substrate should host a twin, and
the literature answers that question far less often than it answers the
question of what a twin is for.

Two families of answer dominate. The first is the integrated platform: managed
cloud services such as Azure Digital Twins and open frameworks such as Eclipse
Ditto provide a graph of twin objects, an ingestion pipeline, and a query
surface, but they assume a cloud-centric deployment and treat the twin as a
document or a node in a managed graph rather than as a live computational
entity~\cite{ditto,azuredt}. The second is the actor framework, in which each
twin is modelled as an actor---an independent unit of state and behaviour that
communicates only by messages. This lineage is the more faithful one, because a
population of physical assets, each evolving on its own and interacting through
well-defined couplings, is almost a textbook instance of the actor model of
Hewitt and of Agha~\cite{hewitt1973,agha1986}. Microsoft's Orleans and its
virtual-actor abstraction made this practical at cloud scale, and its ideas
have been carried onto the \beam in the form of Erleans~\cite{orleans,erleans}.
Yet the actor framework is usually retrofitted onto a runtime---the CLR, the
JVM---whose scheduler, memory model, and failure semantics were designed for
something else.

This paper takes the actor reading seriously and follows it to its conclusion.
If a twin is most naturally an actor, then the right substrate is the runtime
for which the actor model is not a library but the native execution model. That
runtime is the \beam, the virtual machine underlying Erlang and the Open
Telecom Platform~\cite{armstrong2003,cesarini2016}. Our central claim is that
the design space the \emph{Handbook of Digital Twins} enumerates---the levels
at which a twin can sit, the maturity ladder from mirror through shadow to
twin, the topology binding twin to asset, the synchronization regime, the
decision loop, and the classes of user a twin must serve---is not an
inconvenient list of features to be implemented on top of a substrate but a
set of requirements that \otp already satisfies through processes, supervision
trees, transparent message passing, distribution, and hot code
loading~\cite{handbook,hamzaoui2024,otpdesign}. Building a \dt framework in pure
Erlang is therefore less an act of construction than one of naming: most of the
hard parts are already mechanisms in the runtime.

We make three contributions. First, we read the Handbook's vendor-neutral
taxonomy as a specification and establish a point-by-point correspondence
between its axes and \beam mechanisms, arguing that the fit is structural rather
than incidental. Second, from that correspondence we derive an actor-oriented
framework whose core abstraction is a twin as a supervised state machine
carrying an explicit observed-versus-modelled state split, with asset topology
expressed directly as a supervision hierarchy and with transport, addressing,
history, and time isolated behind small behaviours. Third, we identify the
edge-to-cloud span as the property that distinguishes a \beam-native framework
from both platforms and ported actor frameworks: the same twin module runs on a
microcontroller under AtomVM and in a cloud cluster with identical semantics,
and a running model can be replaced in place. We ground the design in a
water-distribution case study and are candid about where the approach is weak.

The remainder of the paper is organized as follows. Section~\ref{sec:space}
reconstructs the \dt design space from the Handbook. Section~\ref{sec:why}
argues for the actor model and for the \beam specifically.
Section~\ref{sec:arch} presents the framework architecture.
Section~\ref{sec:edge} develops the edge-to-cloud argument.
Section~\ref{sec:case} describes the water-distribution case study.
Section~\ref{sec:related} situates the work, and Section~\ref{sec:disc}
discusses limitations and future directions before
Section~\ref{sec:conc} concludes.

\section{The Digital-Twin Design Space}
\label{sec:space}
To argue that a runtime fits the \dt problem, one must first fix what the
problem is. We reconstruct it from the Handbook, which is useful precisely
because it is a consensus document assembled from more than a hundred
contributors rather than the position of a single vendor or
school~\cite{handbook}.

\subsection{From the Apollo Twin to the Mirrored Spaces Model}
The genealogy of the idea is instructive. During the Apollo programme NASA
maintained physical duplicates of its spacecraft on the ground, using them for
training and for rehearsing responses to in-flight contingencies; every
high-fidelity prototype that reproduces operating conditions is in this sense a
twin~\cite{agrawal2024}. Because building a physical duplicate of every asset is
impractical, the twinning idea was moved into software. Grieves introduced the
notion of virtually twinning a physical object in a product-lifecycle course in
the early 2000s, naming it the Mirrored Spaces Model and describing it as three
parts: a physical entity, a digital entity, and the data connection that links
their states bidirectionally~\cite{grieves2017,agrawal2024}. The label
``digital twin'' was attached to this construct around 2011, and NASA gave the
first widely cited definition in 2012, casting the twin as an integrated,
multi-physics, multi-scale, probabilistic simulation that uses the best
available models and sensor history to mirror the life of its flying
counterpart~\cite{glaessgen2012}. Two features of this lineage matter for what
follows. The twin is constituted not by the digital artefact alone but by the
coupling of physical and digital together with the data flow between them, and
the binding is bidirectional in intent from the outset.

\subsection{Two Questions, and an Asset's Logical Counterpart}
Agrawal and Fischer reduce the definitional confusion to two questions that
must be answered afresh for every application: in a given twin, what exactly is
\emph{digital}, and what exactly are we \emph{twinning}~\cite{agrawal2024}? The
questions are deceptively simple, and their value is that they force a designer
to name the modelled aspects of the asset and the fidelity at which each is
represented rather than appealing to the word ``twin'' as if it carried a fixed
meaning. The network-engineering treatment in the Handbook sharpens the
vocabulary: a twin couples a \emph{physical object}---a device, product, or
resource---with a \emph{logical object} that virtualizes the physical object's
relevant characteristics in software, and Minerva and colleagues characterize a
twin through properties such as representativeness and contextualization, the
requirement that the logical object be credible enough to stand in for the
physical one within the context of use~\cite{diniz2024,minerva2020}. The
designer's task is to decide which facets of the asset to represent and how
faithfully, not to represent everything.

\subsection{The Maturity Ladder: Mirror, Shadow, Twin}
The most consequential distinction for a framework is the maturity of the
coupling, which the Handbook presents as a progression and which corresponds to
the categorical review of Kritzinger and colleagues~\cite{hamzaoui2024,kritzinger2018}.
At the lowest rung a \emph{digital mirror} is a collection of mathematical
models describing the asset and its behaviour but with no live link to it---a
simulation that runs independently of any particular physical instance. A
\emph{digital shadow} adds a one-way data thread: the digital representation
follows the state of its physical counterpart as that state changes, but
exercises no influence back on it. Only when a feedback path from the digital
side to the physical side is present, whether partial or complete, does the
coupling become a \emph{digital twin} in the strict sense, and richer variants
layer predictive capability to yield a cognitive twin and human collaboration
to yield a collaborative one~\cite{hamzaoui2024}. Wright and Davidson make the
same point from the other direction: the line between a model and a twin is
exactly the live, bidirectional binding to a specific physical
instance~\cite{wright2020}. A framework that aspires to host twins, and not
merely models or shadows, must therefore make the feedback path and the
distinction between observed and modelled state first-class rather than
incidental.

\subsection{Dimensions of a Twin}
Beyond maturity, the Handbook's deployment methodology decomposes a twin along
several further axes that together constitute a design space, and it is this
decomposition that we will later read as a specification~\cite{hamzaoui2024}.
The \emph{level} axis locates a twin in a part--whole hierarchy that runs from a
component, through equipment composed of components, to a system and finally a
system of systems; the same physical artefact may be modelled as an opaque
component in one twin and as a composite with internal structure in another. The
\emph{topology} axis describes how the twin is sited relative to its asset: a
disconnected or pre-twin used for design-time validation before any link
exists, and, once connected, an embedded twin co-located with the asset but
constrained in computation, a disjoint twin linked but physically remote, or a
joint arrangement in which part of the twin is embedded and part remote. The
\emph{synchronization} axis ranges over asynchronous regimes that are
event-based, conditional, or on-demand and synchronous regimes that are
real-time, near-real-time, or periodic, and the rate may itself vary with the
situation. The \emph{decision loop} axis distinguishes an open loop in which a
human acts on the twin's output, a closed loop in which the twin acts on the
asset directly, and a mixed loop between them. Finally, the \emph{users} axis
records that a twin must present suitable interfaces not only to humans but to
devices, to software applications, and---significantly---to other twins. This
last observation, that a twin is routinely a client of other twins, is what
turns a collection of twins into a system and motivates composition as a
first-class concern.

\subsection{Reference Models, Composition, and a Formal View}
These axes sit within the better-known reference models. The Reference
Architecture Model Industrie~4.0 organizes assets along hierarchy, lifecycle,
and layer axes, and Tao and colleagues extend Grieves' original triad of
physical element, virtual element, and data into a five-dimensional model that
adds \emph{services} and \emph{connections} as explicit
dimensions~\cite{hamzaoui2024,tao2019}. The connections dimension---the data
and control links among the physical asset, the virtual model, and the
services around them---is what the framework must realize as transport, and the
services dimension is what it must realize as queryable behaviour. At the level
of method, the Handbook situates twin construction within model-based systems
engineering and observes that the models composing a complex twin differ both
horizontally, by the parts they describe, and vertically, by the concerns or
viewpoints they capture, so that building a twin recapitulates the composition
of the asset itself and demands principled means of integrating heterogeneous
models~\cite{kovalyov2024}. Finally, the Handbook offers a formal lens that we
will reuse. Sulema and colleagues model a twin as a temporally linked sequence
of discrete states, in which the asset occupies exactly one state at each
instant and that state is keyed by time and determined by the multimodal data
available at that instant~\cite{sulema2024}. Writing $T$ for an ordered set of
time keys and $s_t$ for the state at key $t \in T$, a twin is in this view a
time-indexed trajectory
\begin{equation}
\label{eq:tlm}
\Sigma = \{\, (t,\, s_t) \;:\; t \in T \,\},
\end{equation}
together with the machinery that produces, stores, and queries it. We adopt
\eqref{eq:tlm} as the abstract object our framework must instantiate per asset,
and we will refine $s_t$ into observed and modelled parts in
Section~\ref{sec:arch}.

\section{Why the Actor Model, and Why the \beam}
\label{sec:why}
The design space of Section~\ref{sec:space} reads, line by line, as a list of
properties. We now argue that the actor model supplies those properties and
that the \beam supplies the actor model in a form mature enough to build on.

\subsection{A Twin Is an Actor}
An actor, in the formulation of Hewitt and the semantics of Agha, is an
autonomous entity with private state that interacts with other actors only by
asynchronous messages, processing one message at a time and in response
updating its state, sending further messages, and creating new
actors~\cite{hewitt1973,agha1986}. Set this beside the twin of
Section~\ref{sec:space}. A twin owns the state of one asset and no other; it
must not share that state, because representativeness is a property of one
physical instance. It evolves in response to inputs---telemetry from its asset,
commands from a user, ticks of a clock---which is exactly message handling. It
must address and be addressed by other twins, which is the users axis naming
``another twin'' as a peer. A population of assets becomes a population of
actors, their physical couplings become message channels, and the part--whole
level axis becomes nesting of actors within actors. The correspondence is not
an analogy reached for after the fact; it is the same structure described in two
vocabularies. The temporally linked trajectory of \eqref{eq:tlm} is then simply
the history of one actor's successive states.

\subsection{Why Pure Erlang/\otp, and Not an Actor Library}
If the actor model is the right model, one might still implement it as a library
on any runtime, and Orleans on the CLR and Akka on the JVM do exactly
that~\cite{orleans}. The argument for the \beam is that it provides as runtime
guarantees several things those libraries must approximate, and that the gap
matters for twins specifically. \beam processes are not operating-system
threads and not library objects scheduled cooperatively; they are extremely
lightweight, preemptively scheduled units of which a single node sustains
millions, each with its own heap and its own garbage collection, so that one
twin's allocation and collection cannot pause another~\cite{cesarini2016}. This
per-process isolation is the runtime expression of the twin's requirement that
its state be private and its evolution independent. Failure is contained by the
same isolation: a fault in one process does not corrupt another's memory, and
the supervisor idiom turns the handling of failure into a declared part of the
program structure rather than defensive code scattered through
it~\cite{armstrong2003,otpdesign}. Crucially for an unattended twin mirroring a
running asset, supervision means that the failure of a twin is followed by its
automatic restart into a known state---self-healing is the default rather than
an added feature. The \beam further supports replacing the code of a running
system without stopping it, so that the model a twin uses can be updated while
the twin continues to track its asset, and it makes message passing between
processes location-transparent, so that whether two twins run on one node or on
two communicating nodes is a deployment fact and not a difference in the
program. Erlang and \otp were built to keep telephone switches running for
years through faults and upgrades, and the properties that purpose
demanded---isolation, supervision, live upgrade, distribution---are the
properties a long-lived twin of a long-lived asset demands as
well~\cite{armstrong2003,cesarini2016}.

\subsection{The Correspondence, Axis by Axis}
The fit can be made precise against Section~\ref{sec:space}. The level axis,
from component to system of systems, is the supervision hierarchy: a leaf twin
is a worker process, an equipment or system twin is a supervisor over its
constituent twins, and a system-of-systems twin is a supervisor over those, so
that the part--whole structure of the asset and the process tree of the program
are the same tree. The maturity ladder is realized by what flows along the edges
of that tree and back to the asset: a process that only ingests telemetry is a
shadow, and one whose handler emits a command to its asset's adapter closes the
loop into a twin. The topology axis is a deployment choice over where processes
run, with an embedded twin a process on a constrained node and a disjoint twin a
process on a remote node, communicating identically either way because message
passing is transparent. The synchronization axis is the process's choice of when
to act: a twin that reacts to incoming telemetry messages is event-based, while
one that uses the state machine's timeouts to act at fixed intervals is
periodic, and a twin can do both. The decision loop axis is whether and how a
handler issues actuation, an open loop surfacing output for a human and a closed
loop sending a command to the asset. The users axis is the set of senders a twin
will accept messages from---human-facing front ends, device adapters, other
applications, and peer twins addressed through the registry. Even the
heterogeneous, horizontally and vertically differing models of the model-based
view~\cite{kovalyov2024} have a natural home: distinct concerns can be distinct
processes, composed by the same supervision and messaging that compose
everything else. None of these mappings requires a mechanism the runtime lacks;
each names a use of a mechanism it already has.

\section{Framework Architecture}
\label{sec:arch}
We now derive the framework. The guiding decision is the division between what
the framework owns and what a twin author supplies: everything generic---
lifecycle, supervision, addressing, transport, history, and time---is lifted
into reusable infrastructure, and the author writes only domain logic against a
behaviour. A \emph{behaviour}, in \otp, is an interface together with a generic
implementation of the common control flow, so that the author provides a set of
callback functions and inherits the rest; the framework is, in essence, a small
collection of such behaviours.

\subsection{The \texttt{digital\_twin} Behaviour}
The core of the framework is a behaviour we call \texttt{digital\_twin}, layered
over the \otp \texttt{gen\_statem} state-machine engine rather than the simpler
\texttt{gen\_server}. The choice is deliberate. Physical assets have operational
modes---a pump is stopped, starting, running, or faulted; a valve is open,
transitioning, or closed---and a state machine expresses mode-dependent
behaviour and mode transitions directly, while its built-in state timeouts give
a twin a cost-free notion of staleness: a twin that has heard nothing from its
asset within a configured interval transitions itself into a \texttt{stale} mode
and is no longer mistaken for a live mirror. A purely continuous asset whose
state is a single scalar collapses to a one-mode machine, so nothing is lost by
defaulting to the more expressive engine and specializing downward. The author
implements a small set of callbacks: an \texttt{init} that establishes the
twin's initial state and configuration, a \texttt{handle\_observation} invoked
when telemetry arrives, a \texttt{handle\_command} invoked when actuation is
requested, a \texttt{model\_step} that advances the twin's internal model, and a
\texttt{handle\_query} that answers questions about the twin's state. The
framework supplies the process, its placement in the supervision tree, its
registration for addressing, the wiring to adapters, and the recording of
history, so that the author's code concerns only the asset's behaviour.

\subsection{The Observed--Modelled Split and the Divergence}
The single design decision that separates this framework's twins from mere
shadows is to make the state $s_t$ of \eqref{eq:tlm} carry two components
explicitly. The twin holds an \emph{observed} state $s_t^{\mathrm{obs}}$, the
most recent telemetry from the asset, and a \emph{modelled} state
$s_t^{\mathrm{mod}}$, what the twin's own model predicts should be the case. The
\texttt{handle\_observation} callback updates the former; the
\texttt{model\_step} callback advances the latter according to the asset's
dynamics and any commands in effect, so that, writing $u_t$ for the inputs
active over the interval ending at $t$,
\begin{equation}
\label{eq:model}
s_t^{\mathrm{mod}} = f\!\left(s_{t-1}^{\mathrm{mod}},\, u_t\right),
\end{equation}
with $f$ the author's \texttt{model\_step}. The framework then computes the
divergence between observation and model under an author-supplied metric $d$,
\begin{equation}
\label{eq:div}
\delta_t = d\!\left(s_t^{\mathrm{obs}},\, s_t^{\mathrm{mod}}\right),
\end{equation}
and exposes it as part of every twin's state. This divergence is what gives a
twin a reason to compute rather than merely to store, and it is the quantity in
which the phenomena of interest appear: a sensor fault, a modelling error, and a
genuine process anomaly all manifest as $\delta_t$ crossing a threshold, after
which the closed-loop variants may feed the discrepancy back into actuation. In
the vocabulary of Section~\ref{sec:space}, \eqref{eq:div} operationalizes the
line between a shadow and a twin and instruments the maturity ladder directly:
a twin that only maintains $s_t^{\mathrm{obs}}$ is a shadow, one that also
maintains $s_t^{\mathrm{mod}}$ and reports $\delta_t$ is a twin, and one that
acts on $\delta_t$ is a closed-loop twin. The construction is faithful to the
temporally linked model of \eqref{eq:tlm}: each recorded state is keyed by time
and now resolves into observed and modelled parts whose relationship is the
twin's principal output.

\subsection{Adapters and the Connections Dimension}
A twin must not know how its telemetry arrives or how its commands are
delivered, because the same asset may be reached over different transports in
different deployments and because a twin should be testable without any
transport at all. The framework therefore isolates transport behind a second
behaviour, \texttt{dt\_adapter}, whose implementations translate between the
wire---an industrial protocol such as MQTT or OPC~UA, a serial link to a
microcontroller, or a file---and the twin's observation and command messages.
This layer is the connections dimension of the five-dimensional model and the
device-communication entity of the manufacturing twin standard rendered as
code~\cite{tao2019,iso23247}. The most useful adapter for development is one
that drives no hardware at all: a simulator adapter that replays the output of a
physics simulation as though it were live telemetry, which lets the entire
framework be exercised, and a twin's model validated against ground truth,
before any device exists. Whether such an adapter should bridge to a live
co-simulation through a port or replay logged simulator output for
deterministic, repeatable tests is a genuine choice; the deterministic
replay path is the one to build first, because reproducibility is worth more
than liveness during development.

\subsection{Composite Twins and the Topology of the Asset}
The level axis of Section~\ref{sec:space} becomes architecture through
composition. A real asset network---a water distribution system of tanks,
pumps, pipes, and junctions---is a graph, and that graph maps onto a supervision
hierarchy: leaf twins represent single assets, and a \emph{composite} twin owns
a slice of the topology, supervises the twins beneath it, and coordinates them
so that, for instance, a pump twin's computed outflow becomes the inflow input
of the downstream tank twin at each step. Composites nest, a subsystem twin
supervising several composites and a plant twin supervising the subsystems, and
this nesting is exactly the aggregate twin of the manufacturing
standard~\cite{iso23247}. One sub-decision is genuinely open and worth flagging
rather than hiding: a composite may step its children explicitly by message, or
it may merely supervise while leaf twins exchange values directly along declared
edges of the topology. Explicit composite-driven stepping is the more
conservative starting point, because it keeps the advance of logical time
across the network tractable and centralizes the ordering of updates, whereas
peer-to-peer propagation is more decoupled but makes the global time-step harder
to reason about. The framework fixes the mechanism---supervision and
messaging---and leaves the coordination policy to the composite author.

\subsection{Pluggable Registry, Store, and Clock}
Three further concerns are deliberately not fixed, because their right
implementation depends on where the framework runs, and each is hidden behind a
small behaviour so that the choice is a configuration rather than a rewrite. The
\emph{registry} resolves a twin's logical identity to its process; a default
backed by the \beam's built-in process registry suffices on a single node, while
a distributed registry such as Horde supports addressing and migration of twins
across a cluster~\cite{horde,partisan}. The \emph{store} records the temporally
linked history of \eqref{eq:tlm}; a bounded in-memory ring buffer is appropriate
on a constrained device that cannot retain unbounded history, while a real
time-series store is appropriate in the cloud. The \emph{clock} supplies the
twin's notion of time; a wall-clock drives live operation, whereas a logical
clock drives replay and fast-forward over recorded history, which is what makes
the synchronization axis---real-time, periodic, on-demand---a property the
operator selects rather than one the framework hard-codes. Pluggability here is
not gold-plating; it is the mechanism by which one codebase spans the topology
and synchronization axes of the design space.

\subsection{Lifecycle: Resident and Passivated Twins}
A final decision concerns whether twins are always-resident processes or are
activated on demand and passivated when idle, in the manner of the virtual
actors that Orleans introduced and Erleans brought to the
\beam~\cite{orleans,erleans}. The two are not in opposition. A twin that mirrors
a currently running asset should remain resident, because it must react to
telemetry at any moment and because the cost of repeatedly reconstructing its
state would be wasteful; a twin of a dormant asset, or a historical twin
reconstructed for analysis of a past period, is better passivated, materialized
from the store when queried and released when idle. The framework's position is
therefore that residency is itself a property of the twin's mode, with live
twins resident and dormant or historical twins virtual, and the lifecycle
behaviour makes this selectable. This is the point at which the design touches
the classic-versus-virtual-actor question directly, and the resolution we adopt
is the pragmatic one of letting the asset's liveness, rather than a global
policy, decide.

\section{The Edge-to-Cloud Continuum}
\label{sec:edge}
The property that most distinguishes a \beam-native framework from the
alternatives is that it spans the topology axis of Section~\ref{sec:space} in a
single runtime and a single codebase, and this deserves separate treatment
because it is the crux of the argument.

The \beam is no longer confined to the server. AtomVM is an implementation of
the \beam small enough to run on microcontrollers such as the ESP32, the GRiSP
platform brings Erlang to bare-metal embedded hardware, and the Nerves project
packages \beam applications as the entire software image of an embedded
device~\cite{atomvm,grisp,nerves}. Because all of these execute \beam bytecode
and communicate by the same message passing, a twin written against the
\texttt{digital\_twin} behaviour can run as a thin leaf on a sensor-bearing
microcontroller, close to the asset and within the embedded topology of the
design space, and the very same module can run as a composite in a cloud
cluster, with the two communicating as transparently as two processes on one
node. The embedded, disjoint, and joint topologies of Section~\ref{sec:space}
become, on the \beam, a question of which node a process is placed on, decided at
deployment and not at design. A platform twin cannot offer this, because it
assumes the cloud; a ported actor framework cannot offer it, because its runtime
does not descend to the microcontroller. The single-codebase span is therefore
not a convenience but the realization of an axis of the standard taxonomy that
the alternatives realize only by splitting the system in two.

Two further runtime properties compound the advantage. Live code replacement
means a twin's model---the $f$ of \eqref{eq:model}---can be improved and loaded
into a running twin without stopping it, so the twin need not lose synchrony
with an asset that cannot itself be paused, a consideration that is acute
precisely for the long-lived infrastructure assets that twins are most valuable
for. And supervision means that the failure of an edge twin, a tier where faults
are common, is met by automatic restart into a known state rather than by
silent divergence. The combination---descend to the device, distribute to the
cloud, upgrade in place, recover by default---is available together only because
all four are properties of one runtime.

\section{Case Study: A Water-Distribution Twin}
\label{sec:case}
We ground the architecture in the domain of water distribution, which exercises
every axis of the design space and for which a standard benchmark exists. The
C-Town network, derived from a real medium-sized distribution system and
disseminated through the Battle of the Attack Detection Algorithms, comprises
tanks, pumps, valves, and pipes organized into district metered areas, and has
become a common substrate for studies of monitoring and anomaly
detection~\cite{ctown}. Its dynamics can be generated by established
tools---EPANET for hydraulic simulation and the Water Network Tool for
Resilience built upon it---which makes it an ideal source of ground-truth
telemetry for a simulator adapter~\cite{epanet,wntr}.

The canonical instance we use to exercise the framework is the smallest
composite that is not trivial: a pump feeding a tank. The pump twin and the tank
twin are leaf processes under a composite that owns this fragment of topology.
At each logical step the composite advances its children, the pump twin's
\texttt{model\_step} computes an outflow from its mode and setpoint, and that
outflow is delivered as the inflow input to the tank twin, whose own
\texttt{model\_step} integrates it against demand to update the predicted level,
realizing \eqref{eq:model} across two coupled twins. The simulator adapter
replays EPANET-generated telemetry for the same fragment as observations, so
each twin maintains the observed level alongside its modelled level and reports
the divergence of \eqref{eq:div}; a model-in-the-loop check then reduces to
asking whether each twin's modelled trajectory tracks the simulator's
ground-truth trajectory within tolerance, and an injected discrepancy---a
simulated leak absent from the model---appears as a growing $\delta_t$ that a
closed-loop variant could act upon. Because the adapter drives the system from
logged simulation rather than hardware, the entire case study is deterministic
and repeatable, and because the twins are ordinary supervised processes, the
failure of the pump twin mid-run is followed by its restart and re-registration
without disturbing the tank twin. The design described here is consistent with,
and motivated by, a working Erlang prototype of a twin-per-asset system over the
same benchmark; the present paper abstracts that prototype into the framework
whose generality it was built to demonstrate.

\section{Related Work}
\label{sec:related}
The framework sits at the intersection of three bodies of work. The first is the
\dt platform. Azure Digital Twins offers a managed graph of twin instances
described by a modelling language, and Eclipse Ditto provides an open
abstraction for device twins with state synchronization and a query
interface~\cite{azuredt,ditto}. These are mature and operationally capable, but
they are cloud-resident and treat the twin as a managed object rather than a
live computational actor, and neither descends to the embedded topology that the
design space includes. The second body of work is the actor and virtual-actor
runtime. Orleans established the virtual-actor abstraction and demonstrated it at
scale, Akka brought actors to the JVM, and Dapr generalized the pattern as a
sidecar; on the \beam, Erleans ports the virtual-actor idea, Horde provides
distributed process registration and supervision, and Partisan re-engineers the
distribution layer for larger clusters~\cite{orleans,erleans,horde,partisan}.
Our contribution relative to this lineage is not a new actor mechanism but the
observation, developed in Section~\ref{sec:why}, that the twin's requirements
align with the actor model so closely that the native \beam mechanisms suffice,
together with the edge-to-cloud span that the ported frameworks do not provide.
The third body of work is the standards and conceptual literature, including the
manufacturing twin framework of ISO~23247, the concepts and terminology of
ISO/IEC~30173, and the surveys and conceptual models that the Handbook
collects~\cite{iso23247,iso30173,handbook,fuller2020,minerva2020}; we use this
literature as the specification against which the framework is designed rather
than as a competitor. To our knowledge, a digital-twin framework conceived as a
direct realization of this design space in pure Erlang/\otp, and spanning the
edge-to-cloud continuum on a single runtime, has not previously been articulated.

\section{Discussion and Limitations}
\label{sec:disc}
The argument of this paper is that a runtime can fit a problem structurally, and
that the \beam fits the \dt problem because the actor model the runtime
implements natively is the model the twin demands. Honesty requires marking
where the fit is partial. The clearest limitation is numerical: the \beam is
optimized for concurrency, message passing, and fault tolerance, not for the
dense floating-point computation that high-fidelity physical models, finite
element analyses, and large-scale optimization require, and a twin whose model
is computationally heavy will need to delegate that computation through native
implemented functions or to an external solver rather than express it in Erlang.
This is less a refutation than a boundary: the framework's contribution is the
orchestration, lifecycle, and composition of twins, and heavy numerics belong
behind the \texttt{model\_step} callback as a called service. A second
limitation is ecosystem maturity. The platforms enjoy extensive tooling for
modelling languages, visualization, and integration that a research framework
does not, and although the \beam's primitives are mature, twin-specific
libraries on top of them are not, so early adopters carry more of the burden
themselves. A third concerns the open design questions named in
Section~\ref{sec:arch}: whether composites step children explicitly or leaves
propagate peer-to-peer, whether the simulator adapter co-simulates live or
replays logs, and what policy governs passivation are decisions whose best
answers likely depend on the domain and which a single framework should
parameterize rather than legislate.

Several directions follow. The temporally linked formalism of \eqref{eq:tlm}
and the message protocols among composite and leaf twins are amenable to formal
specification and verification, which the \beam community has pursued for
protocol-bearing systems and which would let one prove properties of a twin
network rather than test for them. The edge claim of Section~\ref{sec:edge}
invites empirical substantiation: a quantitative study of a twin running under
AtomVM on representative microcontroller hardware, measuring the footprint and
latency of the \texttt{digital\_twin} behaviour at the edge, would convert the
architectural argument into measured evidence. And because twins of
infrastructure are attractive targets, the security of the twin network---the
authentication of adapters, the integrity of the observation and command
channels, and the containment of a compromised twin by the same isolation that
contains a failed one---is a natural and pressing extension of the present work.

\section{Conclusion}
\label{sec:conc}
The digital twin's definitional fragmentation has obscured a structural fact:
the design space the field has converged on, as codified in the \emph{Handbook
of Digital Twins}, is a description of autonomous, stateful, hierarchically
composed, fault-prone, long-lived, and variably connected entities, and that is
a description of actors. We have argued that the Erlang/\otp runtime, the \beam,
supplies the actor model not as a library but as its native execution model,
and that each axis of the \dt design space---level, maturity, topology,
synchronization, decision loop, and users---maps onto a mechanism the runtime
already provides. From that correspondence we derived a framework in pure
Erlang whose twins are supervised state machines carrying an explicit
observed-versus-modelled split, whose divergence operationalizes the
shadow-to-twin maturity ladder, whose asset topology is a supervision hierarchy,
and whose transport, addressing, history, and time are pluggable behaviours; and
we identified the edge-to-cloud span, in which one twin module runs from
microcontroller to cluster with live upgrade and default recovery, as the
property that distinguishes a \beam-native framework from both platforms and
ported actor frameworks. Building a twin framework on the \beam is, in the end,
less an act of construction than of recognition, because the hardest
requirements are already mechanisms in the runtime. The water-distribution case
study shows the recognition is workable; the limitations show where it must
borrow; and the structural fit suggests the recognition is worth pursuing.

\begin{thebibliography}{00}

\bibitem{handbook} Z. Lyu, Ed., \emph{Handbook of Digital Twins}. Boca Raton, FL, USA: CRC Press, Taylor \& Francis, 2024.

\bibitem{tao2019} F. Tao, H. Zhang, A. Liu, and A. Y. C. Nee, ``Digital twin in industry: State-of-the-art,'' \emph{IEEE Trans. Ind. Informat.}, vol.~15, no.~4, pp.~2405--2415, 2019.

\bibitem{fuller2020} A. Fuller, Z. Fan, C. Day, and C. Barlow, ``Digital twin: Enabling technologies, challenges and open research,'' \emph{IEEE Access}, vol.~8, pp.~108952--108971, 2020.

\bibitem{agrawal2024} A. Agrawal and M. Fischer, ``What is digital and what are we twinning?: A conceptual model to make sense of digital twins,'' in \emph{Handbook of Digital Twins}, Z. Lyu, Ed. Boca Raton, FL, USA: CRC Press, 2024, ch.~2, pp.~13--29.

\bibitem{hewitt1973} C. Hewitt, P. Bishop, and R. Steiger, ``A universal modular ACTOR formalism for artificial intelligence,'' in \emph{Proc. 3rd Int. Joint Conf. Artif. Intell. (IJCAI)}, 1973, pp.~235--245.

\bibitem{agha1986} G. Agha, \emph{Actors: A Model of Concurrent Computation in Distributed Systems}. Cambridge, MA, USA: MIT Press, 1986.

\bibitem{orleans} P. A. Bernstein, S. Bykov, A. Geller, G. Kliot, and J. Thelin, ``Orleans: Distributed virtual actors for programmability and scalability,'' Microsoft Research, Redmond, WA, USA, Tech. Rep. MSR-TR-2014-41, 2014.

\bibitem{erleans} T. Sloughter \emph{et al.}, ``Erleans: Erlang Orleans-style virtual actors,'' Erleans Project. [Online]. Available: \url{https://github.com/erleans/erleans}

\bibitem{armstrong2003} J. Armstrong, ``Making reliable distributed systems in the presence of software errors,'' Ph.D. dissertation, Dept. Microelectron. Inf. Technol., Royal Inst. Technol. (KTH), Stockholm, Sweden, 2003.

\bibitem{cesarini2016} F. Cesarini and S. Vinoski, \emph{Designing for Scalability with Erlang/OTP: Implementing Robust, Fault-Tolerant Systems}. Sebastopol, CA, USA: O'Reilly Media, 2016.

\bibitem{hamzaoui2024} M. A. Hamzaoui and N. Julien, ``A generic deployment methodology for digital twins -- First building blocks,'' in \emph{Handbook of Digital Twins}, Z. Lyu, Ed. Boca Raton, FL, USA: CRC Press, 2024, ch.~7, pp.~102--121.

\bibitem{otpdesign} Ericsson AB, ``OTP design principles,'' Erlang/OTP System Documentation. [Online]. Available: \url{https://www.erlang.org/doc/design_principles/des_princ.html}

\bibitem{grieves2017} M. Grieves and J. Vickers, ``Digital twin: Mitigating unpredictable, undesirable emergent behavior in complex systems,'' in \emph{Transdisciplinary Perspectives on Complex Systems}, F.-J. Kahlen, S. Flumerfelt, and A. Alves, Eds. Cham, Switzerland: Springer, 2017, pp.~85--113.

\bibitem{glaessgen2012} E. H. Glaessgen and D. S. Stargel, ``The digital twin paradigm for future NASA and U.S. Air Force vehicles,'' in \emph{Proc. 53rd AIAA/ASME/ASCE/AHS/ASC Struct., Struct. Dyn. Mater. Conf.}, Honolulu, HI, USA, 2012, AIAA 2012-1818.

\bibitem{diniz2024} P. H. Diniz, C. E. Rothenberg, and J. F. de Rezende, ``When digital twin meets network engineering and operations,'' in \emph{Handbook of Digital Twins}, Z. Lyu, Ed. Boca Raton, FL, USA: CRC Press, 2024, ch.~3, pp.~30--50.

\bibitem{minerva2020} R. Minerva, G. M. Lee, and N. Crespi, ``Digital twin in the IoT context: A survey on technical features, scenarios, and architectural models,'' \emph{Proc. IEEE}, vol.~108, no.~10, pp.~1785--1824, 2020.

\bibitem{kritzinger2018} W. Kritzinger, M. Karner, G. Traar, J. Henjes, and W. Sihn, ``Digital twin in manufacturing: A categorical literature review and classification,'' \emph{IFAC-PapersOnLine}, vol.~51, no.~11, pp.~1016--1022, 2018.

\bibitem{wright2020} L. Wright and S. Davidson, ``How to tell the difference between a model and a digital twin,'' \emph{Adv. Model. Simul. Eng. Sci.}, vol.~7, art.~13, 2020.

\bibitem{kovalyov2024} S. P. Kovalyov, ``Key technologies of digital twins: A model-based perspective,'' in \emph{Handbook of Digital Twins}, Z. Lyu, Ed. Boca Raton, FL, USA: CRC Press, 2024, ch.~6, pp.~85--101.

\bibitem{sulema2024} Y. Sulema, A. Pester, I. Dychka, and O. Sulema, ``Digital twin model formal specification and software design,'' in \emph{Handbook of Digital Twins}, Z. Lyu, Ed. Boca Raton, FL, USA: CRC Press, 2024, ch.~12, pp.~203--220.

\bibitem{iso23247} \emph{Automation Systems and Integration -- Digital Twin Framework for Manufacturing -- Part 1: Overview and General Principles}, ISO Standard 23247-1:2021, International Organization for Standardization, Geneva, Switzerland, 2021.

\bibitem{iso30173} \emph{Digital Twin -- Concepts and Terminology}, ISO/IEC Standard 30173:2023, International Organization for Standardization, Geneva, Switzerland, 2023.

\bibitem{horde} M. Schaefermeier \emph{et al.}, ``Horde: Distributed, fault-tolerant process registry and dynamic supervisor for Elixir,'' Horde Project. [Online]. Available: \url{https://github.com/derekkraan/horde}

\bibitem{partisan} C. Meiklejohn, H. Miller, and P. Alvaro, ``Partisan: Scaling the distributed actor runtime,'' in \emph{Proc. USENIX Annu. Tech. Conf. (USENIX ATC)}, 2019, pp.~63--76.

\bibitem{atomvm} D. Mazza \emph{et al.}, ``AtomVM: The Erlang virtual machine for IoT devices,'' AtomVM Project. [Online]. Available: \url{https://www.atomvm.net/}

\bibitem{grisp} Stritzinger GmbH, ``GRiSP: Bare-metal Erlang for embedded systems.'' [Online]. Available: \url{https://www.grisp.org/}

\bibitem{nerves} Nerves Project Authors, ``Nerves: Craft and deploy bulletproof embedded software in Elixir.'' [Online]. Available: \url{https://nerves-project.org/}

\bibitem{azuredt} Microsoft, ``Azure Digital Twins documentation.'' [Online]. Available: \url{https://learn.microsoft.com/azure/digital-twins/}

\bibitem{ditto} Eclipse Foundation, ``Eclipse Ditto: Digital twins for IoT.'' [Online]. Available: \url{https://eclipse.dev/ditto/}

\bibitem{ctown} R. Taormina \emph{et al.}, ``Battle of the attack detection algorithms: Disclosing cyber attacks on water distribution networks,'' \emph{J. Water Resour. Plan. Manage.}, vol.~144, no.~8, art.~04018048, 2018.

\bibitem{epanet} L. A. Rossman, \emph{EPANET 2 Users Manual}, U.S. Environmental Protection Agency, Cincinnati, OH, USA, Rep. EPA/600/R-00/057, 2000.

\bibitem{wntr} K. A. Klise, M. Bynum, D. Moriarty, and R. Murray, ``A software framework for assessing the resilience of drinking water systems to disasters with an example earthquake case study,'' \emph{Environ. Model. Softw.}, vol.~95, pp.~420--431, 2017.

\end{thebibliography}

\end{document}
